\section{Implementation}
The database is implemented on \textbf{Oracle Database} using the \textbf{Object-Relational} model.
Each entity is defined as an Oracle \texttt{OBJECT TYPE} and stored in an \texttt{OF} table.
References between objects are modelled with \texttt{REF} attributes.
The N:M relationship between \texttt{Team} and \texttt{Depot} is implemented as a standard relational
join table (\texttt{Team\_Depot}) since it carries no additional attributes.

\subsection{Types Definition}

\begin{lstlisting}
CREATE OR REPLACE TYPE TeamTY AS OBJECT (
    Code                VARCHAR2(20),
    Name                VARCHAR2(100),
    Max_Members         INTEGER,
    Installation_Number INTEGER
)
\end{lstlisting}

\begin{lstlisting}
CREATE OR REPLACE TYPE CustomerTY AS OBJECT (
    ID      VARCHAR2(20),
    Name    VARCHAR2(100),
    Type    VARCHAR2(50),
    Email   VARCHAR2(100),
    Phone   VARCHAR2(20),
    Address VARCHAR2(200)
)
\end{lstlisting}

\begin{lstlisting}
CREATE OR REPLACE TYPE DepotTY AS OBJECT (
    ID                  VARCHAR2(20),
    Name                VARCHAR2(100),
    City                VARCHAR2(100),
    Province            VARCHAR2(100),
    Address             VARCHAR2(200),
    Region              VARCHAR2(100),
    Number_Of_Employees INTEGER
)
\end{lstlisting}

\begin{lstlisting}
CREATE OR REPLACE TYPE MunicipalityTY AS OBJECT (
    ID    VARCHAR2(20),
    Name  VARCHAR2(100),
    Depot REF DepotTY
)
\end{lstlisting}

\begin{lstlisting}
CREATE OR REPLACE TYPE MemberTY AS OBJECT (
    ID      VARCHAR2(20),
    Name    VARCHAR2(100),
    Surname VARCHAR2(100),
    Phone   VARCHAR2(20),
    Email   VARCHAR2(100),
    Team    REF TeamTY
)
\end{lstlisting}

\begin{lstlisting}
CREATE OR REPLACE TYPE Event_LocationTY AS OBJECT (
    ID                  VARCHAR2(20),
    Street              VARCHAR2(200),
    House_Number        VARCHAR2(20),
    Postal_Code         VARCHAR2(10),
    City                VARCHAR2(100),
    Province            VARCHAR2(100),
    Setup_Time_Estimate INTEGER,
    Equipment_Capacity  INTEGER,
    Customer            REF CustomerTY
)
\end{lstlisting}

\begin{lstlisting}
CREATE OR REPLACE TYPE BookingTY AS OBJECT (
    ID             VARCHAR2(20),
    Type           VARCHAR2(50),
    Booking_Date   DATE,
    Duration       INTEGER,
    Cost           NUMBER,
    Booking_Method VARCHAR2(50),
    Contract_Type  VARCHAR2(50),
    Team           REF TeamTY,
    Customer       REF CustomerTY,
    Location       REF Event_LocationTY
)
\end{lstlisting}

\subsection{Tables Definition}

\begin{lstlisting}
CREATE TABLE Team OF TeamTY (
    Code PRIMARY KEY,
    Name NOT NULL,
    Max_Members NOT NULL,
    CONSTRAINT chk_max_members CHECK (Max_Members > 0),
    Installation_Number NOT NULL,
    CONSTRAINT chk_installation_number CHECK (Installation_Number >= 0)
)
\end{lstlisting}

\begin{lstlisting}
CREATE TABLE Customer OF CustomerTY (
    ID PRIMARY KEY,
    Name NOT NULL,
    CONSTRAINT chk_customer_type CHECK (Type IN ('Individual', 'Company')),
    Type NOT NULL,
    Email NOT NULL
)
\end{lstlisting}

\begin{lstlisting}
CREATE TABLE Depot OF DepotTY (
    ID PRIMARY KEY,
    Name NOT NULL,
    City NOT NULL,
    Address NOT NULL,
    Number_Of_Employees NOT NULL,
    CONSTRAINT chk_num_employees CHECK (Number_Of_Employees >= 0)
)
\end{lstlisting}

\begin{lstlisting}
CREATE TABLE Municipality OF MunicipalityTY (
    ID PRIMARY KEY,
    Name NOT NULL UNIQUE,
    Depot NOT NULL REFERENCES Depot ON DELETE CASCADE
)
\end{lstlisting}

\begin{lstlisting}
CREATE TABLE Member OF MemberTY (
    ID PRIMARY KEY,
    Name NOT NULL,
    Surname NOT NULL,
    Team NOT NULL REFERENCES Team ON DELETE CASCADE
)
\end{lstlisting}

\begin{lstlisting}
CREATE TABLE Event_Location OF Event_LocationTY (
    ID PRIMARY KEY,
    Street NOT NULL,
    City NOT NULL,
    Setup_Time_Estimate NOT NULL,
    CONSTRAINT chk_setup_time CHECK (Setup_Time_Estimate > 0),
    Equipment_Capacity NOT NULL,
    CONSTRAINT chk_equip_capacity CHECK (Equipment_Capacity > 0),
    Customer NOT NULL REFERENCES Customer ON DELETE CASCADE
)
\end{lstlisting}

\begin{lstlisting}
CREATE TABLE Booking OF BookingTY (
    ID PRIMARY KEY,
    CONSTRAINT chk_booking_type   CHECK (Type IN ('recurring','one-time')),
    CONSTRAINT chk_booking_method CHECK (Booking_Method IN
        ('phone','email','website','postal_mail')),
    CONSTRAINT chk_contract_type  CHECK (Contract_Type IN
        ('one-time','seasonal','promotional')),
    CONSTRAINT chk_duration CHECK (Duration > 0),
    CONSTRAINT chk_cost     CHECK (Cost >= 0),
    Type NOT NULL, Booking_Date NOT NULL,
    Duration NOT NULL, Cost NOT NULL,
    Booking_Method NOT NULL, Contract_Type NOT NULL,
    Team     NOT NULL REFERENCES Team,
    Customer NOT NULL REFERENCES Customer,
    Location NOT NULL REFERENCES Event_Location
)
\end{lstlisting}

\begin{lstlisting}
-- N:M join table (Team <-> Depot, BELONG relationship)
CREATE TABLE Team_Depot (
    Team_Code VARCHAR2(20) NOT NULL,
    Depot_ID  VARCHAR2(20) NOT NULL,
    CONSTRAINT pk_team_depot PRIMARY KEY (Team_Code, Depot_ID),
    CONSTRAINT fk_td_team  FOREIGN KEY (Team_Code) REFERENCES Team(Code)
        ON DELETE CASCADE,
    CONSTRAINT fk_td_depot FOREIGN KEY (Depot_ID)  REFERENCES Depot(ID)
        ON DELETE CASCADE
)
\end{lstlisting}

\subsection{Population Procedure}

\begin{lstlisting}
CREATE OR REPLACE PROCEDURE PopulateDatabase(
    p_num_customers IN NUMBER,
    p_num_bookings  IN NUMBER
) IS
BEGIN
    -- 20 Teams
    FOR i IN 1..20 LOOP
        INSERT INTO Team VALUES (TeamTY(
            'Team' || TO_CHAR(i),
            'TeamName' || TO_CHAR(i),
            ROUND(DBMS_RANDOM.VALUE(5, 15)), 0));
    END LOOP;

    -- 5 Members per team
    FOR team_id IN 1..20 LOOP
        FOR i IN 1..5 LOOP
            DECLARE v_mid VARCHAR2(20) :=
                'MB' || TO_CHAR((team_id-1)*5+i); BEGIN
            INSERT INTO Member VALUES (MemberTY(
                v_mid, 'Name'||v_mid, 'Surname'||v_mid,
                '33'||TO_CHAR(ROUND(DBMS_RANDOM.VALUE(1000000,9999999))),
                'member'||v_mid||'@sagefit.com',
                (SELECT REF(t) FROM Team t WHERE t.Code =
                    'Team'||TO_CHAR(team_id))));
            END;
        END LOOP;
    END LOOP;

    -- 10 Depots, 3 Municipalities each
    FOR i IN 1..10 LOOP
        INSERT INTO Depot VALUES (DepotTY(
            'Depot'||TO_CHAR(i), 'DepotName'||TO_CHAR(i),
            'City'||TO_CHAR(i), 'Province'||TO_CHAR(i),
            'Address'||TO_CHAR(i),
            'Region'||TO_CHAR(CEIL(i/2)),
            ROUND(DBMS_RANDOM.VALUE(10,50))));
        FOR j IN 1..3 LOOP
            INSERT INTO Municipality VALUES (MunicipalityTY(
                'Mun'||TO_CHAR((i-1)*3+j),
                'Municipality'||TO_CHAR((i-1)*3+j),
                (SELECT REF(d) FROM Depot d WHERE d.ID='Depot'||TO_CHAR(i))));
        END LOOP;
    END LOOP;

    -- Team_Depot assignments
    FOR i IN 1..20 LOOP
        INSERT INTO Team_Depot VALUES (
            'Team'||TO_CHAR(i), 'Depot'||TO_CHAR(MOD(i-1,10)+1));
    END LOOP;

    -- Customers and Event_Locations
    FOR i IN 1..p_num_customers LOOP
        INSERT INTO Customer VALUES (CustomerTY(
            'Cust'||TO_CHAR(i), 'CustomerName'||TO_CHAR(i),
            CASE WHEN MOD(i,2)=0 THEN 'Individual' ELSE 'Company' END,
            'cust'||TO_CHAR(i)||'@example.com',
            '33'||TO_CHAR(ROUND(DBMS_RANDOM.VALUE(1000000,9999999))),
            'Via Roma '||TO_CHAR(i)));
        INSERT INTO Event_Location VALUES (Event_LocationTY(
            'Loc'||TO_CHAR(i), 'Via Sport '||TO_CHAR(i), TO_CHAR(i),
            TO_CHAR(70000+MOD(i,100)), 'City'||TO_CHAR(MOD(i,10)+1),
            'Province'||TO_CHAR(MOD(i,5)+1),
            ROUND(DBMS_RANDOM.VALUE(30,180)),
            ROUND(DBMS_RANDOM.VALUE(50,500)),
            (SELECT REF(c) FROM Customer c WHERE c.ID='Cust'||TO_CHAR(i))));
    END LOOP;

    -- Bookings
    FOR i IN 1..p_num_bookings LOOP
        DECLARE
            v_cid VARCHAR2(20) := 'Cust'||TO_CHAR(
                ROUND(DBMS_RANDOM.VALUE(1,p_num_customers)));
            v_lid VARCHAR2(20) := 'Loc'||TO_CHAR(
                ROUND(DBMS_RANDOM.VALUE(1,p_num_customers)));
            v_tc  VARCHAR2(20) := 'Team'||TO_CHAR(
                ROUND(DBMS_RANDOM.VALUE(1,20)));
        BEGIN
            INSERT INTO Booking VALUES (BookingTY(
                'Book'||TO_CHAR(i),
                CASE WHEN MOD(i,2)=0 THEN 'recurring' ELSE 'one-time' END,
                SYSDATE+ROUND(DBMS_RANDOM.VALUE(-60,60)),
                ROUND(DBMS_RANDOM.VALUE(1,8)),
                ROUND(DBMS_RANDOM.VALUE(100,2000)),
                CASE MOD(i,4) WHEN 0 THEN 'phone' WHEN 1 THEN 'email'
                    WHEN 2 THEN 'website' WHEN 3 THEN 'postal_mail' END,
                CASE MOD(i,3) WHEN 0 THEN 'one-time'
                    WHEN 1 THEN 'seasonal' WHEN 2 THEN 'promotional' END,
                (SELECT REF(t) FROM Team t WHERE t.Code=v_tc),
                (SELECT REF(c) FROM Customer c WHERE c.ID=v_cid),
                (SELECT REF(l) FROM Event_Location l WHERE l.ID=v_lid)));
        END;
    END LOOP;
END;
\end{lstlisting}

\subsection{Triggers}

\begin{lstlisting}
-- Trigger 1: Enforce team member capacity
CREATE OR REPLACE TRIGGER trg_member_capacity
BEFORE INSERT ON Member
FOR EACH ROW
DECLARE
    v_current_count NUMBER;
    v_max_members   NUMBER;
    v_team_code     VARCHAR2(20);
BEGIN
    v_team_code := DEREF(:NEW.Team).Code;
    SELECT COUNT(*) INTO v_current_count
      FROM Member WHERE DEREF(Team).Code = v_team_code;
    SELECT t.Max_Members INTO v_max_members
      FROM Team t WHERE t.Code = v_team_code;
    IF v_current_count >= v_max_members THEN
        RAISE_APPLICATION_ERROR(-20040,
            'Error: Team ' || v_team_code ||
            ' has reached its maximum capacity of ' ||
            v_max_members || ' members.');
    END IF;
END;
\end{lstlisting}

This trigger enforces Business Rule BR9: before inserting a new member, it counts the current members
in the target team and compares with \texttt{Max\_Members}. If the limit is already reached, the insert is blocked.

\begin{lstlisting}
-- Trigger 2: Increment Installation_Number on booking insert
CREATE OR REPLACE TRIGGER trg_update_installation_count
AFTER INSERT ON Booking
FOR EACH ROW
DECLARE
    v_team_code VARCHAR2(20);
BEGIN
    v_team_code := DEREF(:NEW.Team).Code;
    UPDATE Team
       SET Installation_Number = Installation_Number + 1
     WHERE Code = v_team_code;
END;
\end{lstlisting}

This trigger maintains the redundant \texttt{Installation\_Number} attribute automatically (BR12).
Every time a booking is inserted, the counter of the responsible team is incremented by one,
avoiding costly aggregation queries at read time (see Section~2 redundancy analysis).

\begin{lstlisting}
-- Trigger 3: Reject booking if Equipment_Capacity = 0
CREATE OR REPLACE TRIGGER trg_check_location_capacity
BEFORE INSERT ON Booking
FOR EACH ROW
DECLARE
    v_capacity NUMBER;
BEGIN
    SELECT l.Equipment_Capacity INTO v_capacity
      FROM Event_Location l
     WHERE REF(l) = :NEW.Location;
    IF v_capacity <= 0 THEN
        RAISE_APPLICATION_ERROR(-20050,
            'Error: Event Location has no available equipment capacity.');
    END IF;
END;
\end{lstlisting}

This trigger enforces Business Rule BR11: a booking cannot be placed at a location with zero equipment capacity.

\subsection{Operations}

\begin{lstlisting}
-- Procedure 1: Register a new customer
CREATE OR REPLACE PROCEDURE proc_register_customer (
    p_id      IN VARCHAR2,
    p_name    IN VARCHAR2,
    p_type    IN VARCHAR2,  -- 'Individual' or 'Company'
    p_email   IN VARCHAR2,
    p_phone   IN VARCHAR2,
    p_address IN VARCHAR2
) AS
BEGIN
    INSERT INTO Customer
    VALUES (CustomerTY(p_id, p_name, p_type, p_email, p_phone, p_address));
    COMMIT;
EXCEPTION
    WHEN OTHERS THEN
        RAISE_APPLICATION_ERROR(-20001,
            'Error in proc_register_customer: ' || SQLERRM);
END;
\end{lstlisting}

\begin{lstlisting}
-- Procedure 2: Add a new booking
CREATE OR REPLACE PROCEDURE proc_add_booking (
    p_id             IN VARCHAR2,
    p_type           IN VARCHAR2,
    p_date           IN DATE,
    p_duration       IN INTEGER,
    p_cost           IN NUMBER,
    p_method         IN VARCHAR2,
    p_contract_type  IN VARCHAR2,
    p_team_code      IN VARCHAR2,
    p_customer_id    IN VARCHAR2,
    p_location_id    IN VARCHAR2
) AS
    v_count_team     NUMBER;
    v_count_customer NUMBER;
    v_count_location NUMBER;
BEGIN
    SELECT COUNT(*) INTO v_count_team     FROM Team          WHERE Code=p_team_code;
    SELECT COUNT(*) INTO v_count_customer FROM Customer       WHERE ID=p_customer_id;
    SELECT COUNT(*) INTO v_count_location FROM Event_Location WHERE ID=p_location_id;
    IF v_count_team=0 THEN
        RAISE_APPLICATION_ERROR(-20010,'Error: Team not found!');
    END IF;
    IF v_count_customer=0 THEN
        RAISE_APPLICATION_ERROR(-20011,'Error: Customer not found!');
    END IF;
    IF v_count_location=0 THEN
        RAISE_APPLICATION_ERROR(-20012,'Error: Event Location not found!');
    END IF;
    INSERT INTO Booking VALUES (BookingTY(
        p_id, p_type, p_date, p_duration, p_cost, p_method, p_contract_type,
        (SELECT REF(t) FROM Team           t WHERE t.Code=p_team_code),
        (SELECT REF(c) FROM Customer       c WHERE c.ID=p_customer_id),
        (SELECT REF(l) FROM Event_Location l WHERE l.ID=p_location_id)));
    COMMIT;
EXCEPTION
    WHEN OTHERS THEN
        RAISE_APPLICATION_ERROR(-20013,
            'Error in proc_add_booking: ' || SQLERRM);
END;
\end{lstlisting}

\begin{lstlisting}
-- Procedure 3: Add a new team
CREATE OR REPLACE PROCEDURE proc_add_team (
    p_code        IN VARCHAR2,
    p_name        IN VARCHAR2,
    p_max_members IN INTEGER
) AS
BEGIN
    INSERT INTO Team VALUES (TeamTY(p_code, p_name, p_max_members, 0));
    COMMIT;
EXCEPTION
    WHEN OTHERS THEN
        RAISE_APPLICATION_ERROR(-20020,
            'Error in proc_add_team: ' || SQLERRM);
END;
\end{lstlisting}

\begin{lstlisting}
-- Procedure 4: Add a new member to a team
CREATE OR REPLACE PROCEDURE proc_add_member (
    p_id      IN VARCHAR2,
    p_name    IN VARCHAR2,
    p_surname IN VARCHAR2,
    p_phone   IN VARCHAR2,
    p_email   IN VARCHAR2,
    p_team    IN VARCHAR2
) AS
    v_count_team NUMBER;
BEGIN
    SELECT COUNT(*) INTO v_count_team FROM Team WHERE Code=p_team;
    IF v_count_team=0 THEN
        RAISE_APPLICATION_ERROR(-20030,'Error: Team not found!');
    END IF;
    INSERT INTO Member VALUES (MemberTY(
        p_id, p_name, p_surname, p_phone, p_email,
        (SELECT REF(t) FROM Team t WHERE t.Code=p_team)));
    COMMIT;
EXCEPTION
    WHEN OTHERS THEN
        RAISE_APPLICATION_ERROR(-20031,
            'Error in proc_add_member: ' || SQLERRM);
END;
\end{lstlisting}

\begin{lstlisting}
-- Procedure 5: Get all members of a specific team (cursor output)
CREATE OR REPLACE PROCEDURE proc_get_team_members (
    p_team_code IN  VARCHAR2,
    p_cursor    OUT SYS_REFCURSOR
) AS
BEGIN
    OPEN p_cursor FOR
        SELECT m.ID, m.Name, m.Surname, m.Phone, m.Email
          FROM Member m
         WHERE DEREF(m.Team).Code = p_team_code;
END;
\end{lstlisting}
